\chapter{Artifacts}

Artifacts are files you want to keep after a job finishes.
They are ideal for reports, binaries, and diagnostic outputs.

\section{Uploading artifacts}
\begin{lstlisting}[language=YAML]
- name: Upload test report
  uses: actions/upload-artifact@v4
  with:
    name: junit-report
    path: reports/junit.xml
    retention-days: 7
\end{lstlisting}

\section{Downloading artifacts}
Artifacts can be downloaded by later jobs in the same workflow.

\begin{lstlisting}[language=YAML]
- name: Download report
  uses: actions/download-artifact@v4
  with:
    name: junit-report
\end{lstlisting}

\section{Artifacts from matrices}
Matrix jobs upload multiple artifacts with unique names.
Include the matrix value in the artifact name.

\begin{lstlisting}[language=YAML]
- name: Upload build
  uses: actions/upload-artifact@v4
  with:
    name: build-${{ matrix.os }}
    path: dist/
\end{lstlisting}

\section{Retention and cost}
Set retention to the shortest window that still enables debugging.
Large artifacts increase storage and download time.

\begin{warningbox}{Warning: artifacts are not secrets}
Artifacts are stored by GitHub and accessible to anyone with access to the run.
Do not upload secrets or credentials.
\end{warningbox}

\section{Artifact design heuristics}
Artifacts should be explicit and minimal:
\begin{itemize}[nosep]
  \item name them by purpose and scope
  \item keep retention short for large files
  \item avoid uploading the entire workspace
\end{itemize}

\section{Artifacts vs releases}
Artifacts are temporary and tied to a run.
Releases are durable and public-facing.
Use artifacts for debugging and staging, and releases for distribution.

\section{Compression and size}
Compress large artifacts before upload when possible.
Large files slow down downloads and can degrade developer feedback loops.

\begin{checklistbox}{Checklist: artifacts}
\begin{itemize}[nosep]
  \item Upload only what you need to keep.
  \item Use meaningful names.
  \item Keep retention short for large files.
  \item Never store secrets in artifacts.
\end{itemize}
\end{checklistbox}
